\chapter{Before We Begin}

\chapteropening{I}{magine you inherited a library.} Thousands of texts, in a language you're slowly learning, written by one of the most precise minds in human history. Some of them are simple and stunning --- a conversation by a river, a parable about a raft, a list of things to pay attention to when you sit down and close your eyes. Others are dense, technical, and seem to repeat the same formulas over and over until your eyes glaze.

You know this collection contains something extraordinary. You've tasted it. Some passages have changed the way you see your own mind. But you also know you're missing things --- connections between texts, structures you can't quite see, layers of meaning that the words seem to point at without fully revealing. You feel like someone handed you the blueprints for a cathedral, but you're reading them as a grocery list.

Now imagine that one of the original architect's best students sat down and wrote you a manual. Not another teaching. Not more philosophy. A \textit{manual} --- here is how to read these texts. Here are the sixteen different ways the teacher structures his lessons. Here are the five methods for cracking open any passage and finding what's inside. Here are the eighteen root concepts that everything connects back to. Learn these, and you'll never read a sutta the same way again.

That manual exists. It's called the Nettippakaraṇa --- literally ``The Guide'' --- and it's been sitting in the Pali Canon for over two thousand years.

So why hasn't anyone read it?

Because the only English translation, by the great Bhikkhu Ñāṇamoli, published in 1962, is --- let's be honest --- almost impossible to follow unless you already understand what it's trying to say. Ñāṇamoli was a brilliant translator, and his rendering is faithful. But faithfulness to a text this compressed, this technical, this elliptical, produces something that reads like an instruction manual for a machine you've never seen, written in a dialect you almost-but-don't-quite speak. Scholars cite it. Nobody curls up with it.

\sectionbreak

\section*{What This Book Does Differently}

This is the first volume in a series called \textit{Read the Buddha's Original Words in Modern English}. The idea behind the series is simple but radical: take the actual canonical texts of early Buddhism --- not summaries, not commentaries, not books \textit{about} the texts, but the texts themselves --- and translate them so they read the way they were meant to land. Clearly. Directly. Like someone talking to you.

The ancient Pali texts were not written to be obscure. They were composed to be understood by real people --- monks, nuns, laypeople, kings, farmers, courtesans --- in real-time, often in a single sitting. The difficulty we experience today is not in the ideas. The ideas are startlingly clear. The difficulty is in the \textit{translation layer}: the Pali compressed into English that preserves the grammatical structure of the original at the cost of its readability.

This translation takes a different approach. We go behind the curtain. Instead of preserving the compression of the original --- where a single Pali compound might pack an entire paragraph of meaning --- we unpack it. We explain what the author is actually doing. We translate not just the words but the \textit{thinking}.

\section*{How This Book Works}

The main text of each chapter is written in plain, modern English. No Pali appears in the body text. When Mahākaccāyana makes a move, we explain the move. When he gives an example, we walk through why that example works. When the original text uses a technical term, we translate it on the spot, in context, so you never have to flip to a glossary mid-sentence (though we have one at the back if you want it).

The Netti contains beautiful summary verses --- ancient poetry that encapsulates each concept in memorable form. We present these as modern English poems that capture the meaning and, where possible, the rhythm. These are not word-for-word cribs --- they're genuine attempts to make the poetry land in English the way it lands in Pali: as something you could actually memorize and carry with you.

Then there are the Notes.

The Notes at the end of each chapter are where the scholarly work lives. This is where we cite the original Pali, explain exactly how we got from the ancient phrasing to our English rendering, discuss variant readings, point to parallel passages in other texts, argue with previous translations, and generally do everything that would make the main text unreadable if we put it there. If you're a Pali student, the Notes are where you'll live. If you're not, you can ignore them entirely and lose nothing essential.

The result is a book that works on two levels simultaneously. Read only the main text and you get a clean, engaging, complete guide to how the Buddha structured his teachings. Read the Notes and you get a scholarly apparatus that traces every interpretive decision back to the Pali source. Both layers are the real book. Choose your adventure.

\section*{Who This Book Is For}

Anyone who has ever wanted to read the Pali Canon with deeper understanding. That includes:

\begin{itemize}[nosep]
\item Meditators who want to understand the texts they're practicing from
\item Students of Buddhism who find the suttas repetitive and want to understand why they're structured that way (spoiler: it's not repetition --- it's architecture)
\item Pali students who want to see how the tradition itself interpreted its own technical vocabulary
\item Scholars who want a readable companion to Ñāṇamoli's literal translation
\item Anyone who's curious what a ``user's manual for enlightenment literature'' actually looks like
\end{itemize}

You don't need to know Pali. You don't need to have meditated. You don't need to be a Buddhist. You just need to be curious about what might be the most sophisticated reading method ever devised for a body of spiritual literature --- and willing to discover that a text from the third century BCE has something genuinely useful to teach you about how to read.

\vspace{2\baselineskip}
\hfill Sarasota, March 2026

\hfill L. Lopin

\cleardoublepage
